\chapter{Few bosons and few fermions}
\label{ch:few-particles}
\index{boson}
\index{fermion}

\section{Bosonic occupation space}
For $N$ indistinguishable bosons in $m$ modes, a Fock basis state is an occupation vector
\[
(n_1,\ldots,n_m),
\qquad
n_1+\cdots+n_m=N.
\]
The dimension is
\begin{equation}
D_{\rm boson}(m,N)
=
\binom{N+m-1}{N}.
\label{eq:boson-dim}
\end{equation}
For fixed $N$, this is $O(m^N)$.

\section{Fermionic occupation space}
For $N$ fermions in $m$ modes, each mode is occupied at most once and
\begin{equation}
D_{\rm fermion}(m,N)
=
\binom mN.
\label{eq:fermion-dim}
\end{equation}
For fixed $N$, this is also polynomial in $m$.

\begin{corollary}[Few-particle collapse]
A polynomial-size number-preserving bosonic or fermionic circuit with a fixed number of particles is classically polynomial-time simulable if its initial occupation coordinates, mode-gate updates, and final observable are computable in polynomial time in the occupation basis.
\end{corollary}

\begin{proof}
Use \cref{eq:boson-dim,eq:fermion-dim} and \cref{thm:poly-packet}.
\end{proof}

\begin{figure}[H]
\centering
\includegraphics[width=.82\textwidth]{few_particle_dimensions.pdf}
\caption{Fixed-particle bosonic and fermionic sectors grow polynomially with the number of modes.}
\end{figure}

\section{Gate updates}
A number-preserving one-body unitary acts on mode creation operators by
\[
a_i^\dagger
\longmapsto
\sum_j U_{ji}a_j^\dagger.
\]
A generic dense mode transformation can induce many occupation-basis transitions, but for fixed $N$ the basis dimension remains polynomial. Local hopping and phase gates are sparse. Few-body interactions remain polynomial to apply for fixed $N$, although the exponent can grow with $N$.

\section{Where the theorem stops}
If $N$ scales linearly with $m$, both dimensions become exponential. The fixed-number symmetry still reduces the full Fock space, but it does not prove polynomial simulation. Boson-sampling-type interference also shows that a fixed total particle number is not enough when the particle count itself scales with input size and the output asks for detailed many-body amplitudes.

\begin{scopebox}
The corollary is a fixed-particle theorem. It is not a claim that general many-boson or many-fermion quantum dynamics are classically easy.
\end{scopebox}
